\section{Workload Guarantees, Reserve Recourse, and Staffing}
\subsection{Why Staffing Should Be Derived from Overload Risk}
A protection portfolio can look strong on a map and still fail because the stations serving that portfolio are overloaded too often. For that reason, the staffing problem is not treated here as a patrol-hours accounting exercise and not as an \(M/M/c\) queueing approximation. It is treated as a workload-guarantee problem in which each station must keep its overload probability below a target.

\subsection{Station Workload After Portfolio Deployment}
Let \(N_{i,t}\) be the number of incidents generated in cell \(i\) on day \(t\), and let \(s(i)\) be the station responsible for first service. Once the portfolio \(a\) is deployed, the daily workload hitting station \(s\) is
\begin{equation}
A_{s,t}
=
\sum_{i\in\mathcal{Z}_s} N_{i,t}\bigl(1-p_i(a)\bigr),
\label{eq:arrival}
\end{equation}
where \(\mathcal{Z}_s\) is the station's service region. Thus protection quality directly reduces workload, which is a more transparent relationship than translating everything through a stylized queue.

\subsection{Backlog Dynamics}
Let \(c_s\) denote the number of two-person teams assigned to station \(s\), measured in daily service units, and let \(B_{s,t}\) denote unfinished work at the end of day \(t\). We model station pressure by the reflected recursion
\begin{equation}
B_{s,t+1}=\max\{0,\;B_{s,t}+A_{s,t}-c_s\}.
\label{eq:erlang}
\end{equation}
This equation is simple but operationally meaningful. A station fails not when its mean load is large in the abstract, but when daily arrivals repeatedly exceed available service capacity and backlog accumulates.

\subsection{Overload Guarantee}
Suppose the workload process \(A_{s,t}\) has mean \(\mu_s\) and sub-Gaussian scale \(\sigma_s\). Then a one-sided concentration bound gives
\begin{equation}
\Pr\!\bigl(A_{s,t}>c_s+b^\star\bigr)
\le
\exp\!\left(
-\frac{(c_s+b^\star-\mu_s)^2}{2\sigma_s^2}
\right),
\label{eq:intercept}
\end{equation}
where \(b^\star\) is the acceptable single-day backlog tolerance. To guarantee
\(\Pr(A_{s,t}>c_s+b^\star)\le\epsilon_s\), it is sufficient to choose
\begin{equation}
c_s
\ge
\mu_s-b^\star+\sqrt{2\sigma_s^2\log(1/\epsilon_s)}.
\label{eq:rangermin}
\end{equation}
This is the new staffing logic. It is a concentration-based workload guarantee, not an average-delay rule.

\subsection{Cell-Level Service Reliability and Protection Index}
Each cell inherits the overload behavior of its serving station. We define
\begin{equation}
g_i = 1-\Pr\!\bigl(A_{s(i),t}>c_{s(i)}+b^\star\bigr),
\label{eq:service}
\end{equation}
and then compute the park-wide protection index by
\begin{equation}
\PI
=
1-
\frac{
\frac{1}{|\Omega|}\sum_{\omega\in\Omega}Y_\omega(a)
}{
\sum_i L_i+\psi\sum_i U_i+\varepsilon
}.
\label{eq:pi}
\end{equation}
Equation~\eqref{eq:pi} measures the fraction of baseline liability and ambiguity resolved by the deployed portfolio.

\subsection{Reserve Recourse Rule}
Mobile reserve is used to relieve overload risk where it matters most. If a temporary reserve package adds \(r_s\) units of effective service to station \(s\), then its marginal benefit is
\begin{equation}
\Delta_s
=
\Pr\!\bigl(A_{s,t}>c_s+b^\star\bigr)
-
\Pr\!\bigl(A_{s,t}>c_s+r_s+b^\star\bigr).
\label{eq:dispatchscore}
\end{equation}
The reserve is assigned to the station maximizing
\[
\mu_s\Delta_s+\vartheta \sigma_s,
\]
which favors stations with both high expected load and high volatility. This is a recourse rule for bad days, not a static hotspot rule.

\begin{figure}[H]
  \centering
  \includegraphics[width=\linewidth]{fig17_feedback_architecture.png}
  \caption{Workload-guarantee architecture linking portfolio deployment, station load, overload probability, and reserve reassignment in one rolling operating loop.}
  \label{fig:feedbackarch}
\end{figure}

\subsection{Daily Service Calendar}
The park still runs staggered patrol departures, continuous sentinel watch, and designated aerial windows. What changes is the scheduling objective. Windows are timed to protect contract cells first and to support the stations whose workload volatility most threatens the service covenant.

\begin{table}[H]
\centering
\caption{Workload-aware reserve and scheduling logic.}
\label{alg:scheduler}
\begin{tabular}{p{0.95\linewidth}}
\toprule
\textbf{Input:} portfolio \(a\), station loads \(A_{s,t}\), team counts \(c_s\), reserve capacity, overload targets \(\epsilon_s\). \\
\textbf{Output:} daily service calendar and reserve assignment. \\
\midrule
1. Compute station workload means \(\mu_s\) and scales \(\sigma_s\) from Equation~\eqref{eq:arrival}. \\
2. Evaluate overload probability using Equation~\eqref{eq:intercept}. \\
3. Rank stations by weighted reserve benefit from Equation~\eqref{eq:dispatchscore}. \\
4. Schedule patrol departures to protect contract cells first. \\
5. Assign reserve to the highest-ranked station or sector. \\
6. Recompute workload guarantees and roll the plan forward. \\
\bottomrule
\end{tabular}
\end{table}

\begin{figure}[H]
  \centering
  \includegraphics[width=\linewidth]{fig12_station_response_map.png}
  \caption{Service-drag landscape on the official basemap, highlighting where the station system is structurally most exposed to overload.}
  \label{fig:respmap}
\end{figure}

\begin{figure}[H]
  \centering
  \includegraphics[width=\linewidth]{fig07_schedule_timeline.png}
  \caption{Workload-aware daily calendar with contract-first patrol departures, persistent watch periods, and reserve windows for volatile stations.}
  \label{fig:schedule}
\end{figure}
\FloatBarrier

\subsection{Baseline Performance}
Under baseline resources, the new contract portfolio materially outperforms uniform deployment. The gains are now interpreted as liability reduction and overload control rather than area coverage.

\begin{table}[H]
  \centering
  \caption{Baseline performance under the ecological-liability contract framework.}
  \label{tab:results}
  \begin{tabular}{lccc}
    \toprule
    Metric & Uniform deployment & Contract portfolio & Change \\
    \midrule
    Protected-liability ratio & 0.65 & 0.88 & +35.4\% \\
    Mean event confirmation & 0.46 & 0.67 & +45.7\% \\
    Station-overload probability & 0.34 & 0.16 & -52.9\% \\
    Unresolved-incident fraction & 0.33 & 0.19 & -42.4\% \\
    Protection index \(\PI\) & 0.57 & 0.76 & +0.19 \\
    \bottomrule
  \end{tabular}
\end{table}

\subsection{Staffing Floor}
The ranger floor is obtained by satisfying Equation~\eqref{eq:rangermin} station by station:
\[
n_{\min}
=
2\sum_s
\left\lceil
\mu_s-b^\star+\sqrt{2\sigma_s^2\log(1/\epsilon_s)}
\right\rceil.
\]
Under baseline assumptions, this lands near 34 field rangers. The binding constraint is no longer route-hours alone; it is the need to keep overload risk below the service covenant on bad days as well as average days.

\subsection{Interpretation}
Two conclusions follow. First, Etosha behaves like a liability ridge, not a uniform corridor. Second, the nonlinear operational failure is overload, not merely undercoverage: once one or two stations lose workload margin, the value of every other asset degrades quickly.

\begin{figure}[H]
  \centering
  \includegraphics[width=\linewidth]{fig08_simulation_results.png}
  \caption{Protection trajectory under the new contract portfolio, showing faster stabilization once liability and overload risk are both controlled.}
  \label{fig:simulation}
\end{figure}

\begin{figure}[H]
  \centering
  \includegraphics[width=\linewidth]{fig09_robustness_distribution.png}
  \caption{Distributional backtest showing that the contract portfolio improves both average protection and the high-loss tail.}
  \label{fig:robustness}
\end{figure}
\FloatBarrier
